\begin{inlinetable}
\centering
\caption{105,000 次适配迭代后的扰动响应增益。数值 0.25、2、16 与 64 表示源图像中位强度（DN），增益为无量纲量，相对增益采用时间中心化。}\label{tab:gains}
\footnotesize
\setlength{\tabcolsep}{4pt}
\begin{tabular}{lrrrr}
\toprule
扰动／推断方式 & R & $G_1$ & $G_2$ & B \\
\midrule
共同 0.25 / T1 & 0.6211 & 0.6059 & 0.6017 & 0.6720 \\
共同 0.25 / T7 & 0.6125 & 0.5982 & 0.5930 & 0.6685 \\
共同 2 / T1 & 0.7967 & 0.8218 & 0.8158 & 0.8311 \\
共同 2 / T7 & 0.7454 & 0.7626 & 0.7577 & 0.7476 \\
共同 16 / T1 & 0.5654 & 0.4591 & 0.4711 & 0.5585 \\
共同 16 / T7 & 0.5702 & 0.4628 & 0.4663 & 0.5723 \\
共同 64 / T1 & 0.1402 & 0.0651 & 0.0750 & 0.1975 \\
共同 64 / T7 & 0.0934 & 0.0763 & 0.0796 & 0.1621 \\
相对 2 / T1 & 0.7887 & 0.7893 & 0.7849 & 0.8501 \\
相对 2 / T7 & 0.6083 & 0.5046 & 0.5242 & 0.5850 \\
\bottomrule
\end{tabular}
\end{inlinetable}
